\documentclass[11pt]{article}
\usepackage{paperreview}
\title{Blind Response Prediction in Reciprocal Resonator Networks\\
A Methods Paper and Prospective RF Test}
\author{Luqman Mahmood}
\date{30 September 2026}
\hypersetup{pdftitle={Blind Response Prediction in Reciprocal Resonator Networks: A Methods Paper and Prospective RF Test},pdfauthor={Luqman Mahmood}}
\begin{document}
\maketitle
\draftstatus{SM4, v1.0.0}{Methods and prospective protocol; no physical RF result.}
\begin{abstract}
Reciprocity and an observable finite-rank modal response can determine a
port response that has not been measured. We set out a prospective test
using a seven-resonator RF network, two reference ports, and a previously
unacquired cross-port pairing. The mathematical completion is established
linear algebra; the experimental question is whether a loaded, imperfect
board and its measurement chain satisfy the required rank and calibration
conditions well enough to make an accurate forecast. Headline unopened
transmission and acquired-but-withheld reflections are scored separately.
A related, but different, synthetic reduced-array study gives sixteen successful
numerical forecasts in the presence of a fixed dynamic instrument;
ordinary rational interpolation agrees to numerical precision. That result
supports a bounded estimator chain, not superiority over conventional
circuit theory. This paper distinguishes ideal graph identities, numerical
evidence, loaded-board assumptions and unperformed physical predictions.
\end{abstract}

\section{What would count as a result?}
The target is a complex response forecast committed before its target
port pairing is acquired, with a stated uncertainty budget and refusal
conditions. A physical pass would validate a bounded predictive protocol
on an RF network.

The conventional comparator is the same reciprocal circuit or modal model
under the same data access. A deliberately under-ranked predictor is useful
as a sensitivity control, but cannot by itself establish an advantage over
a competent conventional method. A fair tie remains a useful qualification
result and must be reported as such. Reduced-port reconstruction of
multiport scattering matrices is established metrology \cite{TippetSpeciale}.
Low-rank skeleton completion and Nystr\"om extension are established
mathematics \cite{Skeleton,Nystrom}; the latter is a related real-symmetric
positive-kernel setting, not an identification with complex reciprocity.
Rational fitting of frequency
responses, including multiport data, is an established baseline
\cite{VectorFit,MatrixVF}.

\section{Completion at an admitted modal rank}
Partition a complex-symmetric cluster response into reference ports $A$
and target ports $B$:
\[
 K=\begin{pmatrix}K_{AA}&K_{AB}\\K_{BA}&K_{BB}\end{pmatrix},
 \qquad K^T=K .
\]
Transpose, not Hermitian conjugation, expresses this reciprocity convention.

\begin{proposition}[Observable-rank completion]\label{prop:completion}
If $\rank K=\rank K_{AA}=r$, then
\begin{equation}\label{eq:completion}
 K_{BB}=K_{BA}K_{AA}^{+}K_{AB},
\end{equation}
where $+$ denotes the Moore--Penrose inverse. The rank condition includes
the required row and column range conditions.
\end{proposition}
\begin{proof}
Since $K_{AA}$ has the full rank of the block matrix, the columns of
$K_{AB}$ lie in its range and the rows of $K_{BA}$ lie in its row space.
Block elimination with $K_{AA}^{+}$ gives a Schur complement
$K_{BB}-K_{BA}K_{AA}^{+}K_{AB}$. Any nonzero rank of that complement
would increase the rank of $K$, contradicting the hypothesis.
\end{proof}
Noise generally destroys exact rank. The actual estimator must fix a
rank-truncation rule and a conditioning cutoff before observing the target.
A graph certificate alone does not prove the loaded response has that rank,
or that the chosen ports see all modes.

\subsection{Why unresolved clusters are the object}
A nearly degenerate pair may span a well-determined subspace even when its
individual residue vectors are not stably separated. We therefore target
band-selected cluster residue matrices, with the band, background subtraction,
normalization and pole/rank admission fixed in advance.
Rank-one completion of a genuine rank-two cluster is a comparator, not an
alternative definition of success. A fitted full-rank background must not
be silently counted as a modal residue or discarded after the target is known.

\subsection{A conditional error budget}
For clarity, choose an invertible $r$-by-$r$ reference block $A$, and write
the prediction as $CA^{-1}B$. If
$\sigma_{\min}(A)\geq s>0$ and errors obey
$\|\Delta A\|\leq\varepsilon_A<s$,
$\|\Delta B\|\leq\varepsilon_B$ and
$\|\Delta C\|\leq\varepsilon_C$, then
\[
 \|\widehat C\widehat A^{-1}\widehat B-CA^{-1}B\|
 \leq
 \frac{\varepsilon_C\|\widehat B\|}{s-\varepsilon_A}
 +\frac{\|C\|\varepsilon_A\|\widehat B\|}{s(s-\varepsilon_A)}
 +\frac{\|C\|\varepsilon_B}{s}.
\]
This follows by adding three differences and using the resolvent identity.
The bound is only as credible as the externally justified error inputs.
It is not a perturbation theorem for a pseudoinverse across a rank change.

\section{Ideal seven-resonator design versus the loaded board}
The ideal comparator is the affine-$E_6$ seven-node star: one center, three
inner nodes and three outer nodes. Its adjacency spectrum is
$-2,-1,0,1,2$ with multiplicities $1,2,1,2,1$.
For equal inductances and compensated diagonal capacitance,
\[
 C=C_{\rm tot}I-C_cA,\qquad
 f_\lambda=\frac{f_0}{\sqrt{1-r\lambda}},\qquad
 r=C_c/C_{\rm tot},\quad
 f_0=\frac1{2\pi\sqrt{LC_{\rm tot}}}.
\]
The design values $L=1\,\mu{\rm H}$, $C_c=68\,{\rm pF}$ and
$C_{\rm tot}=272\,{\rm pF}$ give $r=1/4$ and ideal cluster frequencies
approximately $7.879$, $8.631$, $9.650$, $11.143$ and $13.647$ MHz.
These are design calculations, not measured resonance frequencies.

The proposed board must be modeled with its loads. A probe capacitor in series with a
termination contributes $sC_p/(1+sZ_0C_p)$, not simply $sC_p$.
Loss, unequal inductances, parasitics and asymmetric four-port loading
alter modes, splittings and frequencies. Node-wise tuning does not make
all inductances equal or restore exact ideal symmetry.
The physical protocol must admit the calibrated loaded response, not
require a real board to equal the ideal graph.

Two outer nodes are reference ports $P1,P2$; two inner nodes are
target ports $P3,P4$. The ideal rank-two reference block has condition
number three. That number is an exact ideal-design property, not a
promise about the measured reference block.

\section{Prospective acquisition and scoring protocol}
The thresholds below are design targets; a run-specific protocol must fix
them before acquisition. The supplied protocol annex lists the measurement
conventions, numerical estimator and calibration decisions required before
any new target capture.

\begin{enumerate}
\item Qualify the analyzer, complex direction convention, calibration plane,
loads, component quality and board. Preserve refused admission attempts.
The Touchstone test must include an asymmetric fixture: a reciprocal
fixture alone cannot reveal a transposed off-diagonal writer.
\item Acquire the five reference pairings $P1P2$, $P1P3$, $P1P4$, $P2P3$
and $P2P4$. Fix the admitted bands, loaded fit, rank choices and
uncertainties, and record the prediction file.
\item Commit the prediction before acquiring $P3P4$.
Only $S_{34}$ and $S_{43}$ in that last pairing are genuinely unopened.
The earlier pairings already excite $P3,P4$ and acquire their reflections.
\item Score $K_{34}$ across the five clusters as the primary unopened
prediction. Score the acquired-but-withheld reflections and the aggregate
as separate consistency outcomes.
\item For center, inner and outer-node perturbations, record a new prediction
before each corresponding target pairing, using the same specified regime
and classification rules.
\end{enumerate}

\begin{center}
\begin{tabular}{p{.57\linewidth}p{.28\linewidth}}
\toprule
Outcome & Design relative-$L^2$ bound\\
\midrule
Unopened $K_{34}$ (headline) & $\leq0.10$\\
Withheld self entries & $\leq0.15$\\
Withheld aggregate & $\leq0.12$\\
\bottomrule
\end{tabular}
\end{center}
A withheld failure does not erase headline success, but must remain visible:
the report must state the separate outcome of each metric.
The aggregate and self metrics must use vector norms fixed in the
run-specific protocol, not post-hoc elementwise tolerances.

The perturbation prediction uses the calibrated generalized capacitance
pencil with a specified added $33\,{\rm pF}$ at a named node.
Shifts/splittings must remain in the preselected adjacent-gap regime;
outside it no orbit classification is assigned. The classification margin is
a runner-up-to-best residual ratio of at least two.
The attenuator-ratio check tests frequency-shape consistency only.
An unchanged file or frequency-flat gain can pass it; it does not establish
amplitude linearity.

\subsection{Comparator addition required before a new strong claim}
Alongside the rank-one control, the proposed conventional comparator is
multiport vector fitting followed by reciprocal modal synthesis, using the
same five reference pairings and no target-pair data. Pole/order selection,
shared-pole constraints, residue-rank rule, background, calibration,
regularization and scoring must be fixed before capture. A fit that cannot
identify the unseen coupling from its admitted model must report that
failure, not acquire extra target information.

\section{Related numerical evidence, kept in its own domain}
A separate fixed-instrument synthetic intervention study produced sixteen
numerical forecasts inside their declared radii \cite{RF}. The reduced
arrays are supplied in \texttt{supplement/data/}, and
\texttt{supplement/replay\_ri7.py} reconstructs their prediction bounds.
Here an \emph{instrument} is an additive, frequency-dependent matrix
response held fixed between calibration and target. A \emph{refusal}
means that the calculation withholds a forecast because its rank or
fixed-instrument assumptions fail. Reported normalized
prediction errors were $0.3682$--$1.6977$ parts per million; radii were
$0.01400$--$0.08365$ percent of predicted norm. It also recorded eight
rank refusals and two changing-instrument refusals.
The fixed additive dynamic instrument had a matrix-norm contribution of
$35.11$--$41.60$ percent in that construction; these are not energy fractions.

Ordinary same-data rational interpolation agreed to about
$5.252\times10^{-16}$. Thus the study corroborates a bounded numerical
prediction/error chain under its specified instrument model.
Its $10^{-7}$ complex-entry budget is a declared numerical error model.

The synthetic study's controlled-intervention rank is not automatically the modal rank in
\eqref{eq:completion}. The RF proposal does not reproduce that study merely by
using two reference ports. The two studies motivate related questions
about observability, but their assumptions and credits remain separate.

\section{What a physical pass would and would not establish}
A successful blinded loaded-board forecast would connect a preparation,
a measured response class and a particular unseen response. Useful follow-ups
include a fresh board, altered loading, controlled rank loss and a
predeclared out-of-regime perturbation. Failures can locate a wrong signal
class, incomplete modal observation, unstable extraction or calibration drift;
these causes must be distinguished in reporting the outcome.

A future results paper must include raw complex sweeps, reference-plane
metadata, prediction provenance, every refused or failed attempt and uncertainty
calibration. At present there is no physical result in this manuscript.

\aistatement
\begin{thebibliography}{9}
\bibitem{TippetSpeciale} J. C. Tippet and R. A. Speciale,
\emph{A Rigorous Technique for Measuring the Scattering Matrix of a
Multiport Device with a 2-Port Network Analyzer},
IEEE Trans. Microwave Theory Tech. \textbf{30}, 661--666 (1982).
\url{https://doi.org/10.1109/TMTT.1982.1131118}.
\bibitem{Nystrom} C. K. I. Williams and M. Seeger,
\emph{Using the Nystr\"om Method to Speed Up Kernel Machines},
Advances in Neural Information Processing Systems \textbf{13},
682--688 (2001; conference 2000).
\url{https://proceedings.neurips.cc/paper/2000/hash/19de10adbaa1b2ee13f77f679fa1483a-Abstract.html}.
\bibitem{Skeleton} S. A. Goreinov, E. E. Tyrtyshnikov and N. L. Zamarashkin,
\emph{A theory of pseudoskeleton approximations},
Linear Algebra Appl. \textbf{261}, 1--21 (1997).
\url{https://doi.org/10.1016/S0024-3795(96)00301-1}.
\bibitem{VectorFit} B. Gustavsen and A. Semlyen,
\emph{Rational approximation of frequency domain responses by vector fitting},
IEEE Trans. Power Delivery \textbf{14}, 1052--1061 (1999).
\url{https://doi.org/10.1109/61.772353}.
\bibitem{MatrixVF} Z. Drma\v{c}, S. Gugercin and C. Beattie,
\emph{Vector Fitting for Matrix-valued Rational Approximation},
SIAM J. Sci. Comput. \textbf{37}, A2346--A2379 (2015).
\url{https://doi.org/10.1137/15M1010774}.
\bibitem{RF} L. Mahmood, Reproducibility and prospective-protocol
supplement (2026). See \texttt{supplement/README.md}.
\end{thebibliography}
\end{document}
